A large share of outpatient care in South Asia is delivered not in hospitals but in private
medical chambers, which are one- or two-room practices where a single physician sees patients
in sequence, often with one assistant and no clinical information system. Two routine tasks
dominate the overhead. The first is ordering the queue, typically managed with handwritten
tokens, a paper register, or informal agreement among waiting patients. The second is
recording basic vital signs, which usually happens inside the consultation room after the
patient sits down, so it consumes the one resource the chamber has least of. Measurements of
consultation length in Bangladeshi district and upazila hospitals put the average encounter
at roughly nine minutes~\cite{das2021appointment}, and every minute a thermometer reading and
a pulse count consume within that window is a minute not spent on history taking.

Both tasks are mechanical, and both are amenable to automation. The question is what to automate them
with. The dominant answer in the literature is a networked one. Instrument the patient, stream
readings to a cloud dashboard, and let clinicians reach them through a web or mobile front
end~\cite{irjet2020,rath2021,hrtemp2019}. That approach suits hospitals, and it aligns with the
interoperability goals of the WHO global digital health strategy~\cite{who2021digital}. It is a poor
fit for a two-room chamber, because it adds recurring connectivity cost, makes a consultation depend
on an uplink that may be intermittent, and introduces a failure class the chamber staff cannot
diagnose. That infrastructure assumptions imported from well-resourced settings tend not to survive
contact with developing-region deployments is by now a familiar argument~\cite{brewer2005case}, and a
queue that stalls because a Wi-Fi link dropped is worse than a queue managed on paper.

The complementary work on embedded queue and token systems is offline by construction but scoped to a
different problem. Those systems arbitrate service order at bank counters and ticket windows,
sometimes adding infrastructure such as a GSM link~\cite{qgsm2013} and otherwise confining themselves
to issuing a token number~\cite{eqc2022,arpn2019}. None captures clinical data.

This paper takes the position that the two functions belong in the same device. We present a
stand-alone system, built on one ATmega32 8-bit microcontroller, that guides a patient through
self-registration, acquires body temperature and pulse rate before the patient enters the
consultation room, assigns a queue position, and presents the record to the physician on a
console advanced by a single button. Four requirements shaped it.

\textbf{R1, no network dependency.} Registration, queueing and display must complete with no
uplink, host PC or companion smartphone. \textbf{R2, bounded response.} The physician's control
action must produce a visible result quickly and repeatably. \textbf{R3, operable by
non-technical users.} A first-time patient must register from on-screen prompts alone.
\textbf{R4, commodity parts.} Every component must be locally available at a price a
single-physician practice can justify.

R1 and R4 push toward an 8-bit microcontroller rather than an application processor, and that is what
makes the engineering non-trivial. The ATmega32 offers 32 general-purpose I/O lines, 2~KB of SRAM, and
one 10-bit ADC multiplexed across eight channels~\cite{microchip2020}. From that budget the system
must scan a 32-key matrix, drive two character displays, sample two analog sensors and poll a control
button.

We contribute a complete embedded design integrating FCFS queue arbitration, guided self-registration
with free-text name entry and pre-consultation vital acquisition on one ATmega32, with the full I/O
and memory allocation reported (Sections~\ref{sec:hardware}--\ref{sec:firmware}); a dual-corner
architecture separating patient-side input and sensing from doctor-side review and dequeue control
(Section~\ref{sec:architecture}); the clock constraint we resolved to keep a future serial uplink
viable (Section~\ref{sec:calibration}); and a bench evaluation with an explicit statement of its
limits (Section~\ref{sec:results}).
